%%%%%%%%%%     PREÁMBULO      %%%%%%%%%%%%%%%

%%%%%%%%%   Non modificar este preámbulo, salvo que se sepa ben o que se está a facer   %%%%%%%%%

%%%%%%  Traballando nun Mac, se usas TeXShop, en  Preferencias -> Código -> codificación de  caracteres (UTF-8)  %%%%

\documentclass[11pt,a4paper]{book}

% Formato do documento
\usepackage[papersize={210mm,297mm},lmargin=2.5cm,rmargin=2.5cm,top=3.5cm,bottom=3.5cm]{geometry}   % Marxes
\renewcommand{\baselinestretch}{1.25}                                                               % Fixamos o interliñado
\setlength{\parskip}{8pt}                                                                           % Separación entre parágrafos


\usepackage{amsmath,amsfonts,amssymb,amsthm}
\usepackage{latexsym,pdfsync,xcolor,graphicx}
\usepackage[nottoc]{tocbibind}
\usepackage[T1]{fontenc}
\usepackage[utf8]{inputenc}					
\usepackage[spanish]{babel}			
\usepackage{tikz}
\usepackage{appendix}
\usepackage{hyperref}








% Cabeceira
\usepackage{fancyhdr}
\pagestyle{fancy}
\fancyhf{}
\renewcommand{\chaptermark}[1]{\markboth{\textbf{\thechapter. #1}}{}} % Formato para o capítulo: N. Nome
\renewcommand{\sectionmark}[1]{\markright{\textbf{\thesection. #1}}}  % Formato para a sección: N.M. Nome

\fancyhead[LO]{\rightmark}           % Nas páxinas impares, parte esquerda do encabezado, aparecerá o nome do capítulo
\fancyhead[RE]{\leftmark}            % Nas páxinas pares, parte dereita do encabezado, aparecerá o nombre da sección
\fancyhead[RO,LE]{\textbf{\thepage}} % Números de páxina nas esquinas dos encabezados




%%%   \pagestyle{plain}   Cando os títulos dos capítulos son moi longos, a opción por defecto das cabeceiras pode ser inapropiada;
%%%  nestes e noutros casos semellantes, pódese usar a opción "plain", que suprime as cabeceiras, deixando soamente a numeración das páxinas.


%%%%%%%      NON ESCRIBIR AQUÍ  (COMANDO PORTADA) %%%%%%%%%%%

\newcommand{\portada}[3]{
\thispagestyle{empty}

\definecolor{azulUSC}{RGB}{13,38,120}

\includegraphics[width=13cm]{logomath.pdf}

\vspace*{2cm}

\centerline{{\Large\bf  Traballo Fin de Grao}}

\vspace{2.6cm}

{\center{\Huge\bfseries\color{azulUSC} #1}\par}		

\vspace{1cm}

{\Large

\centerline{#2}		

\vspace{6.4cm}

\centerline{\sf #3}        

\vspace{1.25cm}

\centerline{\sf UNIVERSIDADE DE SANTIAGO DE COMPOSTELA}
}

\enlargethispage{1cm}
\clearpage
}

%%%%%%%      NON ESCRIBIR AQUÍ  (COMANDO PRIMEIRA PÁXINA) %%%%%%%%%%%

\newcommand{\primeira}[3]{
\thispagestyle{empty}
\enlargethispage{1cm}
\begin{center}\Large

{\sf GRAO DE MATEM\'ATICAS}

\bigskip
{\bfseries Traballo Fin de Grao}
\vspace{4cm}

{\bfseries\Huge #1} 		
\vspace{1.2cm}

{#2}		\par		     
\end{center}

\vspace{6.5cm}

\begin{center}
\Large
{#3}	\par			    

\vspace{1cm}
{\sf UNIVERSIDADE DE SANTIAGO DE COMPOSTELA}
\end{center}
\clearpage
}

%%%%%%%    Teoremas   %%%%%%%%%

\newtheorem{theorem}{Teorema}[chapter]
\newtheorem{proposition}[theorem]{Proposición}
\newtheorem{lemma}[theorem]{Lema}
\newtheorem{corollary}[theorem]{Corolario}
\newtheorem{algorithm}[theorem]{Algoritmo}
\newtheorem{axiom}[theorem]{Axioma}
\newtheorem{case}[theorem]{Caso}
\newtheorem{conclusion}[theorem]{Conclusión}
\newtheorem{condition}[theorem]{Condición}
\newtheorem{conjetura}[theorem]{Conjetura}				
\newtheorem{criterion}[theorem]{Criterio}
\newtheorem{ejercicio}[theorem]{Ejercicio}	
			
\theoremstyle{definition}
\newtheorem{definition}[theorem]{Definición}
\newtheorem{ejemplo}[theorem]{Ejemplo}				
\newtheorem{agradecimientos}[theorem]{Agradecimientos}	

\theoremstyle{remark}
\newtheorem{remark}[theorem]{Observación}
\newtheorem{notation}[theorem]{Notación}
\newtheorem{problem}[theorem]{Problema}
\newtheorem{solution}[theorem]{Solución}

\addto\captionsspanish{
	\renewcommand{\contentsname}{\'Indice} 
	\renewcommand\appendixname{Anexo}
	\renewcommand\appendixpagename{Anexos}
}
%%%%%%%%%%%%%%%%%%   Fin do Preámbulo   %%%%%%%%%%%%%%%%




\begin{document}
\frontmatter

%%%%%%%%%%%      PORTADA      %%%%%%%%%%%%

\portada
{Resolución numérica del problema no lineal de mínimos cuadrados.\vspace{5pt}
Aplicaciones a la estimación de parámetros de modelos matemáticos.}		%  escribir aquí o  Título
{Dídac Blanco Morros}		%  escribir aquí o/a  Autor/a
{Curso Acad\'emico}		  	%  escribir aquí a Data de presentación, na forma:  mes, ano

%%%%%%%%     PRIMEIRA PÁXINA    %%%%%%%%%%%%%%%%%%

\mbox{}
\thispagestyle{empty}
\clearpage

\setcounter{page}{1}

\primeira
{Resolución numérica del problema no lineal de mínimos cuadrados. 
Aplicaciones a la estimación de parámetros de modelos matemáticos.}		%  escribir aquí o  Título
{Dídac Blanco Morros}		%  escribir aquí o/a  Autor/a
{Febrero, 2022}		  	%  escribir aquí a Data de presentación, na forma:  mes, ano

\thispagestyle{empty}

%%%%%%%%%%%      PROPOSTA DE TRABALLO      %%%%%%%%%%%%%%%%

\vspace*{.5 cm}

\chapter*{Trabajo propuesto}

\vspace{1 cm}

\begingroup
\renewcommand*{\arraystretch}{2.2}
\begin{tabular}{|l|}
	\hline
	
	{\bf \'Area de Co\~necemento:  \/ }\\ \hline
	\begin{minipage}{11.5cm}
		{\vspace*{.2cm}
			\bf T\'{\i}tulo:   \/ \bf
			\vspace{.2cm}}
	\end{minipage}\\ \hline
	\bf Breve descrici\'on do contido\\ \hline
	\begin{minipage}{11.5cm}
		{\vspace*{.2cm}

		
		\vspace{.2cm}}
	\end{minipage}\\ \hline
	{\bf Recomendaci\'ons}\\ \hline
	\begin{minipage}{11.5cm}
		{\vspace*{.2cm}


		\vspace{.2cm}}
	\end{minipage}\\ \hline
	{\bf Outras observaci\'ons}\\ \hline
	\begin{minipage}{11.5cm}
		{\vspace*{.2cm}


			\vspace{.2cm}}
	\end{minipage}\\ \hline
\end{tabular}
\endgroup

\clearpage

\thispagestyle{empty}

\tableofcontents

\clearpage

\thispagestyle{empty}

\mbox{}

\clearpage

%%%%%%%%%     RESUMO---RESUMEN---ABSTRACT     %%%%%%%%%%%%%

\phantomsection
\addcontentsline{toc}{chapter}{Resumen}		
\chapter*{}

\section*{Resumen}

% escribir aquí resumo en español




\vspace{1.5cm}

\section*{Abstract}



% escribir aquí resumo en inglés



\clearpage

\thispagestyle{empty}

%%%%%%%%%    INTRODUCIÓN     %%%%%%%%%%%%%%%%%%
%%%%%%%%%    De non incluír unha INTRODUCIÓN, suprímese    %%%%%%%%%%

\phantomsection
\chapter*{Introducci\'on}\addcontentsline{toc}{chapter}{Introducci\'on}

\markboth{INTRODUCCI\'ON}{INTRODUCCI\'ON}

%%%%%%%%%%%     fin da INTRODUCIÓN     %%%%%%%%%%%%%

\mainmatter

%%%%%%%%     Póñense os Capítulos    %%%%%%%%%%%%%%%

\chapter{Título del Capítulo 1}



\section{Sección}



\chapter{Título del Capítulo 2}




\chapter{Título del Capítulo 3}




%%%%%%%%     Póñense a bibliografía, Apéndices e o índice alfabético     %%%%%%%%
\appendix
\renewcommand{\thechapter}{\Roman{chapter}}
\chapter{Título del Anexo I}


\chapter{Título del Anexo II}

\backmatter

\begin{thebibliography}{99}

%%%%%%%%----Exemplo de entradas bibliográficas:
%
\bibitem{HiSt} Hilton, P. J. y Stammbach, U. A. (1997). \emph{A course in homological algebra}, 2nd ed., Graduate Texts in Mathematics, 4, Springer-Verlag,  New York.
\bibitem{Lo} Loday, J.-L. (1996). \emph{Künneth-style formula for the homology of Leibniz algebras}, Math. Z., \textbf{221}, 41--47.
%
\end{thebibliography}




\end{document}


